\chapter{Esquemas, conclusiones y cierre de revisión}
\section{Esquemas producidos}
\begin{figure}[htbp]\centering
\includegraphics[width=\linewidth]{figuras/esquema-hidraulico.pdf}
\caption{H-00. Secuencia de aprovechamiento y responsabilidades del tramo derivado.}
\end{figure}
La captación, paso ecológico, purgas y restitución son partes de un mismo aprovechamiento. Las descargas requieren diseño de erosión, calidad y seguridad; no se da por inocua una purga porque vuelve al río.

\begin{figure}[htbp]\centering
\includegraphics[width=\linewidth]{figuras/perfil.pdf}
\caption{H-01. Perfil energético idealizado entre cámara y referencia de turbina. La línea no es levantamiento del terreno.}
\end{figure}
La línea de energía parte del nivel de cámara y desciende con pérdidas modeladas. La línea piezométrica en el interior de la tubería se obtiene restando $v^2/(2g)$ a la energía; la presión local depende además de la cota del eje. No se infiere seguridad frente a vacío a partir del perfil recto simplificado. El máximo transitorio no está dibujado como garantía.

\begin{figure}[htbp]\centering
\includegraphics[width=\linewidth]{figuras/unifilar-funcional.pdf}
\caption{E-01. Unifilar conceptual con barra BT y derivación a servicios propios. Sin ajustes ni secciones aprobadas.}
\end{figure}
E-01 permite revisar responsabilidades y secuencia del sistema. El unifilar definitivo deberá incluir número de polos, conductores, secciones, barras, tierra, neutros, TC/TV, protecciones, relaciones, enclavamientos, grupo vectorial y límites de propiedad. La simbología de construcción se normalizará con las referencias DGE y convenciones del proyecto.

\clearpage
\section{Planos que faltan para un expediente de detalle}
\begin{longtable}{L{1.5cm}L{5.7cm}L{7.2cm}}
\toprule Código & Plano requerido & Entrada que falta\\\midrule
PL-01 & Ubicación, predios y esquema general georreferenciado & Levantamiento y verificación de derechos/afectaciones\\
PL-02 & Perfil hidráulico y topográfico & Eje real, niveles, secciones y restitución\\
PL-03 & Captación y desarenador & Avenidas, sedimentos, geotecnia y diseño estructural\\
PL-04 & Cámara, rejas y aliviadero & Niveles de maniobra, vórtices, excedentes y acceso\\
PL-05 & Tubería y anclajes & Perfil, transitorio, código mecánico y suelos\\
PL-06 & Casa de máquinas y cimentación & Plano de fabricante, izaje, sismo y ventilación\\
PL-07 & Disposición del grupo & Fabricante, cargas dinámicas y mantenimiento\\
PL-08 & Unifilar y lógica de control & Estudio de conexión, fallas y protecciones\\
PL-09 & PAT, neutros y equipotencialidad & Resistividad, corrientes y tiempos de falla\\
PL-10 & Línea MT, perfiles y estructuras & Trazo, cruces, vanos, servidumbre y aislamiento\\
PL-11 & Auxiliares, iluminación y emergencia & Planta y cuadro de cargas validado\\
\bottomrule
\end{longtable}
Las escalas de estos planos se elegirán al disponer de geometría. No se anuncia que se adjuntan planos de detalle inexistentes.

\section{Conclusiones sustentadas}
\begin{enumerate}
\item El borrador no cierra su geometría y energía. Con su caudal, salto neto de 50 m y rendimientos, la salida es \POriginalCincuenta\ kW. Al recalcular su tubería interior de 390 mm y 175 m resulta \PBase\ kW, manteniendo aún una longitud incompatible con las coordenadas.
\item La alternativa prioritaria de 0.420 m$^3$/s, 220 m, diámetro interior 500 mm y salto bruto 52 m produce condicionalmente \PeAlt\ kW en bornes y \PExport\ kW en PCC. Su margen es \MargenExport\ kW. Con 450 mm, la misma carga auxiliar reduce la exportación a \PExportRef\ kW.
\item El caudal mínimo para 150 kW PCC sería \QObjetivo\ l/s con 500 mm bajo estas hipótesis. Una reducción de caudal del 5\% deja la exportación por debajo de la meta. No se acredita potencia firme ni garantía a partir del punto de diseño.
\item Los modos auxiliares explican \PAuxNormal\ kW en generación y \PAuxMax\ kW en mantenimiento; TSA, transformador principal y línea se contabilizan separados. La revisión térmica y de arranque puede exigir cambiar esos valores y recalcular exportación.
\item Las presiones de cribado son \PLento\ bar con cierre de 4 s y \PRapido\ bar para cierre abrupto ideal en 500 mm. Membrana, viga y térmica cuantifican acciones, pero no aprueban material, espesor, apoyos ni PN16 sin código y transitorio ejecutivo.
\item Los datos Q95 y caudal ecológico del borrador no sostienen disponibilidad anual. La nueva alternativa consume todo el saldo ilustrativo de 0.420 m$^3$/s, por lo que la hidrología y derechos deben cerrarse antes de seleccionar potencia.
\item La selección eléctrica requiere garantías en altura, régimen de neutro y puesta a tierra, estudio de conexión y cortocircuito. Se corrigen la descripción Dyn5, la fórmula de altitud y la referencia a Regla 362.
\item La matriz normativa permite revisar qué exige cada norma y qué evidencia falta. Ninguna declaración de cumplimiento integral se infiere de la inclusión de normas en bibliografía.
\item El proyecto queda organizado por capítulos con entradas canónicas y motor reproducible. Las verificaciones numéricas comprueban ecuaciones y coherencia; no reemplazan ensayos, levantamiento ni aprobación técnica.
\end{enumerate}

\section{Orden de cierre recomendado}
Primero se confirma con el docente si Limbani, 150 kW y conexión a red son requisitos o elecciones del alumno. Después se cierra el recurso y la geometría, se compara tecnología con fabricantes y se solicita información de conexión. Con estos datos se desarrollan transitorios, estructuras, protecciones y tierra; recién entonces se fijan metrados y costos definitivos.

La revisión normativa debe incorporar textos consolidados, las reglas complementarias de electrificación rural si la clasificación del sistema las hace aplicables y exigencias específicas de distribuidora. Se deberá obtener la rúbrica del curso para determinar extensión y profundidad final sin alterar la trazabilidad del estudio.

\section{Reproducción y control documental}
El punto de entrada es \texttt{proyecto.yaml}. El constructor lee los archivos JSON de parámetros, normativa y observaciones, calcula resultados y genera tablas y figuras en el directorio de construcción del proyecto. Las fuentes editables del informe están en \texttt{expediente/capitulos/}.

La ejecución produce \texttt{calculos.json}, comparación CSV, observaciones Markdown, huellas SHA-256 y los dos PDF. Se prueban balance energético, Colebrook, geometría imposible, entradas inválidas, maniobra rápida, flujo uniforme y cálculo trifásico. Los archivos técnicos permanecen como revisión R02 hasta validación del docente y de los datos del sitio.
